\section{U-I Characteristics}
\label{sec:UI}

\subsection{Analysis}
In the next part of the experiment, a $4 \times 4 \si{\milli\meter}$ square of polished cristalline silicon and 
a $4 \times 4 \si{\milli\meter}$ square of amorphous silicon are connected to a voltage. First measurements are
made without any influx of light; then a measurement with exactly $ 1 \si{\milli\watt\per\meter}$ of light is made.
The results are shown in figure \ref{fig:UI}.\\
\begin{figure}[h]
  \centering
  \includegraphics[width=\textwidth]{UI.pdf}
  \caption{U-I characteristics of the different pv materials.}
  \label{fig:UI}
\end{figure}
The serial resistance (R) and the parallel resistance (R) can be determined by the dark current measurements.
A fit through the dark current for voltages near 0 is done; resulting in the following values for $R_p$ and $R_s$:
\begin{align*}
  R_{p, crystalline} \approx \dfrac{\del U}{\del I} |_{U_{dark} = 0} = 4.109 \cdot 10^{3} \si{\ohm}\\
  R_{s, crystalline} \approx \dfrac{\del U}{\del I} |_{I_{light} = 0} = 13.66 \si{\ohm}\\
  R_{p, amorphous} \approx \dfrac{\del U}{\del I} |_{U_{dark} = 0} = 7.084 \cdot 10^{7} \si{\ohm}\\
  R_{s, amorphous} \approx \dfrac{\del U}{\del I} |_{I_{light} = 0} = 260.38 \si{\ohm}
\end{align*}
For the serial resistance the derivative at a high voltage in needed. Therefore we use the data of the point
of the $U_ocv$.\\

The Open Circuit Voltage (OCV) is the voltage at which the current is zero. In our case this is at 
$U_{OCV} = 0.5 \si{\volt}$ for crystalline silicon and $U_{OCV} = 0.83 \si{\volt}$ for amorphous silicon.\\
The short circuit current is the current for an external voltage of zero. In our case this is at 
$I_{SC} = -6.709277 \cdotp 10^{-3} \si{\milli\ampere}$ for crystalline silicon and $I_{SC} = -6.97 \cdot 10^{-4} \si{\milli\ampere}$ for amorphous silicon.\\
The Fill Factor (FF) is determined by the following ratio:
\begin{align}
  FF = \dfrac{P_{max}}{I_{SC} \cdot U_{OCV}}
\end{align}
where $P_{max}$ is the maximum power, which is determined maximizing the product of $I(U)$ and $U$. For the probes
we get 
\begin{align}
  FF_{crystalline} = 0.623\\
  FF_{amorphous} = 0.647
\end{align}
Fromthe same data we get the efficiency:
\begin{align}
  \eta_{crystalline} = \dfrac{P_{max}}{P_{light}} = 0.131\\
  \eta_{amorphous} = 0.023
\end{align}

\subsection{Questions}
\textbf{1. How can the serial and parallel resistances be calculated from the dark current characteristics?}\\
From the equivalent Ciruit Model of a Real Solar Cell, discussed in the description of the experiment, we get:
\begin{align*}
  I = I_S (\exp (q(U-IR_S)/AkT) - 1) -I_L + (U -IR_S)/R_P
\end{align*}
For small Voltages, the first term is close to $0$, therefore we can neglect it. Then we get:
\begin{equation*}
  R_p = {\dfrac{\partial U}{\partial I}}|_{V = 0}
\end{equation*}
For high VOltages, the exponential term is much higher, therefor we can neglect the second term. We get
\begin{align*}
  R_s = {\dfrac{\partial U}{\partial I}}|_{high V}
\end{align*}

\textbf{2. The manufacturer of the crystallin solar cell specifies more than $10^{5} \si{\ohm}$ for the parallel 
resistance $R_p$. Why is the measured value much smaller?}\\
Probe insulation, socket leakage, meter input leakage, bench wiring, humidity, or flux/residue on the cell surface 
create extra leakage paths that reduce the measured resistance.

\textbf{3. What are the main sources of error in calculating the efficiency?}\\
The main sources of error in calculating the efficiency come from inaccuracies in measuring the 
I-V characteristics, the incident light power, and the active cell area. In our experimental setup,
the most significant source of error is active cell area, which ia also influenced by the casing. Another 
point of error stems from the fact, that the lamp used for the measurements is, in fact, not a perfect 
representation of the sun.

\textbf{4. Why are the open-circuit voltages and the short-circuit currents of the crystalline and amorphous 
silicon solar cells so different at constant illumination?}\\
Crystalline silicon has a highly ordered structure with long carrier lifetimes and low defect densities, 
leading to higher $V_{oc}$, while amorphous silicon has many defects and localized states that cause strong 
recombination, reducing $V_{oc}$. However, amorphous silicon has a higher optical absorption coefficient, 
so it can generate more current in thin layers, often giving it a higher $I_{sc}$ per thickness even though its 
$V_{oc}$ is lower.